\documentclass[a4paper,11pt]{article}

\usepackage[fontset=fandol]{ctex}
\usepackage{amsmath,amssymb}
\usepackage{graphicx}
\usepackage[a4paper,margin=2.05cm,top=2.25cm,bottom=2.25cm]{geometry}
\usepackage[most]{tcolorbox}
\usepackage{xcolor}
\usepackage{booktabs}
\usepackage{array}
\usepackage{enumitem}
\usepackage{float}
\usepackage{fancyhdr}
\usepackage{hyperref}
\usepackage{listings}
\usepackage{tikz}
\usepackage{etoolbox}

\hypersetup{
    colorlinks=true,
    linkcolor=DeepBlue,
    urlcolor=DeepBlue,
    citecolor=DeepBlue,
    pdftitle={通用面试复盘样例},
    pdfauthor={Codex}
}
\xeCJKsetup{CJKecglue={}}

\definecolor{DeepBlue}{HTML}{1F4E79}
\definecolor{Ink}{HTML}{1F2933}
\definecolor{Slate}{HTML}{52606D}
\definecolor{SoftBlue}{HTML}{EAF4FF}
\definecolor{SoftGreen}{HTML}{EBF8F1}
\definecolor{SoftRed}{HTML}{FFF0F0}
\definecolor{SoftYellow}{HTML}{FFF7DB}
\definecolor{SoftViolet}{HTML}{F3F0FF}
\definecolor{LineGray}{HTML}{D8DEE9}
\definecolor{CodeBack}{HTML}{F7F8FA}

\pagestyle{fancy}
\fancyhf{}
\lhead{\small 通用样例}
\rhead{\small 面试复盘}
\cfoot{\small \thepage}
\renewcommand{\headrulewidth}{0.4pt}
\setlength{\parindent}{0pt}
\setlength{\parskip}{0.55em}
\setlist[itemize]{leftmargin=1.6em,itemsep=0.2em,topsep=0.2em}
\setlist[enumerate]{leftmargin=1.8em,itemsep=0.2em,topsep=0.2em}
\renewcommand{\arraystretch}{1.25}

\lstset{
    basicstyle=\ttfamily\small,
    keywordstyle=\color{DeepBlue}\bfseries,
    stringstyle=\color{red!55!black},
    commentstyle=\color{green!45!black},
    backgroundcolor=\color{CodeBack},
    breaklines=true,
    frame=single,
    rulecolor=\color{LineGray},
    numbers=left,
    numberstyle=\tiny\color{Slate},
    captionpos=b,
    columns=fullflexible,
    keepspaces=true,
    showstringspaces=false,
    tabsize=4
}

\newtcolorbox{verdictbox}[1]{
    enhanced, breakable, sharp corners,
    colback=SoftBlue, colframe=DeepBlue, colbacktitle=DeepBlue,
    coltitle=white, fonttitle=\bfseries, title={#1},
    left=1.1mm,right=1.1mm,top=1mm,bottom=1mm,
    boxrule=0.9pt
}

\newtcolorbox{riskbox}[1]{
    enhanced, breakable, sharp corners,
    colback=SoftRed, colframe=red!60!black, colbacktitle=red!65!black,
    coltitle=white, fonttitle=\bfseries, title={#1},
    left=1.1mm,right=1.1mm,top=1mm,bottom=1mm,
    boxrule=0.9pt
}

\newtcolorbox{betterbox}[1]{
    enhanced, breakable, sharp corners,
    colback=SoftGreen, colframe=green!45!black, colbacktitle=green!45!black,
    coltitle=white, fonttitle=\bfseries, title={#1},
    left=1.1mm,right=1.1mm,top=1mm,bottom=1mm,
    boxrule=0.9pt
}

\newtcolorbox{evidencebox}[1]{
    enhanced, breakable, sharp corners,
    colback=black!2!white, colframe=Slate, colbacktitle=Slate,
    coltitle=white, fonttitle=\bfseries, title={#1},
    left=1.1mm,right=1.1mm,top=1mm,bottom=1mm,
    boxrule=0.8pt
}

\newtcolorbox{drillbox}[1]{
    enhanced, breakable, sharp corners,
    colback=SoftYellow, colframe=yellow!55!black, colbacktitle=yellow!55!black,
    coltitle=black, fonttitle=\bfseries, title={#1},
    left=1.1mm,right=1.1mm,top=1mm,bottom=1mm,
    boxrule=0.8pt
}

\newtcolorbox{codebox}[1]{
    enhanced, breakable, sharp corners,
    colback=CodeBack, colframe=Ink, colbacktitle=Ink,
    coltitle=white, fonttitle=\bfseries, title={#1},
    left=1.1mm,right=1.1mm,top=1mm,bottom=1mm,
    boxrule=0.8pt
}

\newtcolorbox{questioncard}[2][]{
    enhanced, breakable, sharp corners,
    colback=white, colframe=DeepBlue!75!black, colbacktitle=DeepBlue!75!black,
    coltitle=white, fonttitle=\bfseries, title={#2},
    left=1.2mm,right=1.2mm,top=1mm,bottom=1mm,
    boxrule=0.8pt,
    #1
}

\newcommand{\scoretag}[1]{\tikz[baseline=(X.base)] \node[draw=DeepBlue!60,fill=SoftBlue,rounded corners=1.5pt,inner xsep=4pt,inner ysep=2pt](X){\small #1};}
\newcommand{\risktag}[1]{\tikz[baseline=(X.base)] \node[draw=red!55!black,fill=SoftRed,rounded corners=1.5pt,inner xsep=4pt,inner ysep=2pt](X){\small #1};}
\newcommand{\oktag}[1]{\tikz[baseline=(X.base)] \node[draw=green!45!black,fill=SoftGreen,rounded corners=1.5pt,inner xsep=4pt,inner ysep=2pt](X){\small #1};}

\begin{document}

\begin{titlepage}
\thispagestyle{empty}
\centering
{\Large 面试复盘\par}
\vspace{0.6cm}
{\Huge\bfseries 通用面试复盘样例\par}
\vspace{0.35cm}
{\large \par}
\vspace{0.5cm}
{\small 生成时间：YYYY-MM-DD \quad 视频时长：00:08:30 \quad ASR：sample transcript\par}
\vspace{0.7cm}

\begin{verdictbox}{首页摘要：总体结论}
\textbf{总体通过概率}：中等偏上

\textbf{总体结论}：项目主线能讲清，但负责范围、指标口径和代码边界需要更稳。

\textbf{致命风险}：如果被追问缓存一致性或二分 invariant，当前回答有失分风险。
\end{verdictbox}

\vspace{0.35cm}
\begin{minipage}[t]{0.48\textwidth}
\small
\begin{riskbox}{Top 5 扣分点}
\begin{itemize}
\item 项目开头没有先给量化指标。
\item 缓存失效和一致性策略解释偏短。
\item 代码题没有先声明 lower\_bound invariant。
\end{itemize}
\end{riskbox}
\normalsize
\end{minipage}
\hfill
\begin{minipage}[t]{0.48\textwidth}
\small
\begin{betterbox}{Top 5 可复用亮点}
\begin{itemize}
\item 能把业务延迟和缓存方案联系起来。
\item 能承认方案限制并给观测指标。
\item 能在追问下回到数据和边界条件。
\end{itemize}
\end{betterbox}
\normalsize
\end{minipage}

\vspace{0.35cm}
\begin{minipage}[t]{0.48\textwidth}
\small
\begin{drillbox}{高风险技术点}
\begin{itemize}
\item 缓存一致性
\item 超时与取消传播
\item lower\_bound 二分
\end{itemize}
\end{drillbox}
\normalsize
\end{minipage}
\hfill
\begin{minipage}[t]{0.48\textwidth}
\small
\begin{evidencebox}{下一场优先级}
\begin{itemize}
\item 准备 30 秒项目主线。
\item 背熟 lower\_bound 模板。
\item 把实验指标和风险写成固定句。
\end{itemize}
\end{evidencebox}
\normalsize
\end{minipage}

\vfill
\begin{evidencebox}{本地证据说明}
样例使用虚构 transcript 和 review\_plan，不含真实音视频。
\end{evidencebox}
\end{titlepage}

\tableofcontents
\newpage

\section{首页摘要}

\begin{itemize}
\item 这是一个最小样例，用于验证 skill 的 LaTeX 模板、渲染器和验证器。
\item 真实任务中应从本地 transcript、interview\_events 和 evidence frames 构造此 JSON。
\end{itemize}

\section{面试官问题与我的回答}

\begin{questioncard}{1. 缓存系统优化项目 \hfill 01:12--02:40}

\scoretag{passable 3/5} \oktag{key}



\textbf{面试官问题}：你先介绍一下最近做的项目。



\textbf{我的回答}：我最近做的是一个缓存系统优化项目，主要目标是降低核心接口的延迟。我先梳理读路径里哪些数据变化不频繁、访问量又比较高，然后设计缓存 key、TTL 和失效策略，并把命中率、P95/P99 延迟和错误率接到监控里。回答里能体现项目主线，但我没有一开始就把负责模块和最终指标讲清楚。



\textbf{回答来源窗口}：01:15--02:35



\textbf{建议答案}：建议先说：我负责一个缓存优化项目，目标是降低核心接口延迟；我主要做缓存 key 设计、失效策略和监控指标，最后用 P95/P99 延迟、命中率和错误率验证效果。个人贡献和数值必须来自本地转写或用户补充材料。



\textbf{来源}：[L-q01] [REDIS-CACHE]



\textbf{一句评价}：主线能成立，但要先给贡献、指标和限制。

\end{questioncard}

\begin{questioncard}{2. 超时和取消传播 \hfill 03:00--04:20}

\scoretag{risky 3/5} \oktag{key}



\textbf{面试官问题}：如果下游变慢，你怎么避免请求堆积？



\textbf{我的回答}：我当时回答的是入口请求要设置超时，如果下游慢了不能一直等；可以配合重试、限流和降级，避免请求在队列里越堆越多。但我没有把 deadline 向下游传播、取消信号和重试预算讲清楚，也没有强调重试必须满足幂等和总预算限制。



\textbf{回答来源窗口}：03:05--04:10



\textbf{建议答案}：先给原则：入口请求带 deadline，向下游传播 context/cancellation；重试只在幂等且有预算时做，并配合限流、熔断和降级，避免把慢下游打得更慢。



\textbf{来源}：[L-q01] [GO-CONTEXT]



\textbf{一句评价}：有工程直觉，但需要把超时预算和重试风险讲清楚。

\end{questioncard}

\section{重点追问复盘}

\begin{riskbox}{缓存系统优化项目 \hfill 01:12--02:40}

\textbf{追问信号}：验证候选人能否把问题、方案、负责范围和效果压缩成清楚主线。



\textbf{关键扣分点}：主线能成立，但要先给贡献、指标和限制。



\textbf{下次一句话先答}：建议先说：我负责一个缓存优化项目，目标是降低核心接口延迟；我主要做缓存 key 设计、失效策略和监控指标，最后用 P95/P99 延迟、命中率和错误率验证效果。个人贡献和数值必须来自本地转写或用户补充材料。



\textbf{来源}：[L-q01] [REDIS-CACHE]

\end{riskbox}

\begin{riskbox}{超时和取消传播 \hfill 03:00--04:20}

\textbf{追问信号}：验证候选人是否理解超时、取消传播、限流和降级边界。



\textbf{关键扣分点}：有工程直觉，但需要把超时预算和重试风险讲清楚。



\textbf{下次一句话先答}：先给原则：入口请求带 deadline，向下游传播 context/cancellation；重试只在幂等且有预算时做，并配合限流、熔断和降级，避免把慢下游打得更慢。



\textbf{来源}：[L-q01] [GO-CONTEXT]

\end{riskbox}

\section{代码题复盘}

\begin{codebox}{lower\_bound 二分查找 \hfill 05:10--08:00}

\textbf{现场表现}

能写出二分结构，但一开始没有声明搜索区间语义。



\textbf{关键问题}

\begin{itemize}
\item 没有先说返回第一个大于等于 target 的位置。
\item 边界更新解释偏慢。
\end{itemize}



\textbf{标准思路}

维护半开区间 [left, right)，循环条件 left < right；mid 可行时收缩右边界，否则收缩左边界。



\textbf{复杂度}

时间复杂度 O(log n)，空间复杂度 O(1)。



\textbf{下一次口述稿}

我找第一个满足 nums[i] >= target 的位置；右边界用 len(nums)，这样 target 大于所有元素时自然返回 n。



\textbf{参考代码}

\begin{lstlisting}[language=Python]

def lower_bound(nums, target):
    left, right = 0, len(nums)
    while left < right:
        mid = (left + right) // 2
        if nums[mid] >= target:
            right = mid
        else:
            left = mid + 1
    return left

\end{lstlisting}

\textbf{来源}：[L-code01] [PY-BISECT]

\end{codebox}

\section{高风险技术点速记}

\begin{drillbox}{缓存一致性与失效}

\textbf{规则}：先说明数据新鲜度要求，再选择 TTL、主动失效、双写顺序或旁路缓存策略。



\textbf{面试说法}：我会先按业务能接受的数据延迟来选策略：强一致就少缓存或同步失效，弱一致就用 TTL/事件失效，并用命中率、延迟和脏读告警验证。



\textbf{来源}：[REDIS-CACHE]

\end{drillbox}

\begin{drillbox}{超时预算}

\textbf{规则}：每个请求应有总 deadline，子调用共享或切分预算，取消信号向下游传播。



\textbf{面试说法}：我会把超时作为调用链预算，而不是每层独立加一个很长 timeout；超过预算就取消下游并返回降级结果。



\textbf{来源}：[GO-CONTEXT]

\end{drillbox}

\section{参考来源}
\begin{enumerate}
\item [L-q01] \textbf{你先介绍一下最近做的项目。}（local\_event）：interview\_events.json
\item [L-code01] \textbf{代码题：lower\_bound 二分查找}（local\_event）：interview\_events.json
\item [PY-BISECT] \textbf{Python documentation: bisect}（official\_docs）：\url{https://docs.python.org/3/library/bisect.html}；访问日期：YYYY-MM-DD
\item [GO-CONTEXT] \textbf{Go documentation: context package}（official\_docs）：\url{https://pkg.go.dev/context}；访问日期：YYYY-MM-DD
\item [REDIS-CACHE] \textbf{Redis documentation: cache management}（official\_docs）：\url{https://redis.io/docs/latest/develop/reference/client-side-caching/}；访问日期：YYYY-MM-DD
\end{enumerate}

\section{后续巩固资料}
\begin{enumerate}
\item \textbf{二分模板}｜Python documentation: bisect（official\_docs，high，English，short）\par \textbf{为什么看}：用官方 lower\_bound 语义校准返回值、边界和插入点定义。\par \textbf{链接}：\url{https://docs.python.org/3/library/bisect.html}
\item \textbf{缓存系统设计}｜Redis documentation: client-side caching（official\_docs，high，English，short）\par \textbf{为什么看}：用于补齐缓存失效、命中率和一致性风险的面试表达。\par \textbf{链接}：\url{https://redis.io/docs/latest/develop/reference/client-side-caching/}
\item \textbf{超时与取消传播}｜Go documentation: context package（official\_docs，high，English，short）\par \textbf{为什么看}：用于把 deadline、cancellation 和下游调用预算讲成清晰工程机制。\par \textbf{链接}：\url{https://pkg.go.dev/context}
\item \textbf{项目答辩组织}｜来源不足，建议补材料（needs\_user\_material，high，Chinese，short）\par \textbf{为什么看}：个人贡献、实验数值和业务背景只能来自本地视频、简历或用户补充材料。\par \textbf{链接}：来源不足，建议补材料
\item \textbf{实验指标}｜Google SRE Book: Monitoring Distributed Systems（official\_book，medium，English，medium）\par \textbf{为什么看}：用于把延迟、错误率、饱和度和告警指标组织成可防守的工程回答。\par \textbf{链接}：\url{https://sre.google/sre-book/monitoring-distributed-systems/}
\end{enumerate}

\end{document}
